\section{The Headline Result}
\label{sec:disc-headline}

MobileNetV2 achieved higher macro F1 than ResNet-50 under all three imbalance
strategies, and under two of the three the two architectures were statistically
indistinguishable on per-image correctness. A network holding 2.23~million trainable
parameters was never beaten by one holding 23.52~million, on the same data, through the
same pipeline, with the same optimiser and schedule.

The temptation is to read this as MobileNetV2 being the better architecture. That reading
is too strong, and the more defensible one is more interesting: at this dataset size,
capacity beyond MobileNetV2's was not the binding constraint. HAM10000 provides 8{,}012
training images across seven classes, of which the three rarest contribute 467 between
them. ResNet-50 has roughly ten times the parameters to fit from the same evidence, and
the extra capacity had nothing useful to consume. What it did instead is the subject of
Section~\ref{sec:disc-instability}.

That distinction matters for how far the result generalises. It supports the claim that
a lightweight backbone is sufficient for HAM10000-scale dermatoscopic classification. It
does not support a claim that MobileNetV2 would remain competitive given an order of
magnitude more data, and nothing here tests that.

\section{ResNet-50's Optimisation Behaviour}
\label{sec:disc-instability}

The ResNet-50 runs did not descend smoothly. Measuring the standard deviation of the
epoch-to-epoch change in validation loss, averaged within each architecture, gives 0.130
for MobileNetV2 and 0.221 for ResNet-50, a factor of 1.70. E4 produced a single-epoch
increase in validation loss of 1.207 and moved in the wrong direction on 23 of its 49
epoch transitions, close to half.

This is not a defect in the implementation. Every relevant source file was reviewed and
the behaviour is consistent with a known sensitivity: a deep, BatchNorm-heavy network
fine-tuned with Adam at $10^{-3}$, with no warmup, no learning-rate schedule and no
weight decay. Section~\ref{sec:method-regularisation} records those omissions, and they
are the most plausible cause. MobileNetV2 is shallower and its inverted-residual blocks
appear more forgiving of the same configuration.

The learning rate was deliberately not lowered for ResNet-50 alone. Doing so would have
tuned one arm of a comparison whose entire premise is that only the architecture varies,
and the resulting numbers would have answered a different question than the one asked.
The instability is therefore reported as a finding.

It is a finding with practical weight. A network that trains reliably at a shared
learning rate is cheaper to work with than one requiring its own schedule, and that cost
is real even though it appears in no accuracy table. Combined with the efficiency
results, it strengthens rather than complicates the case for the lightweight model.

\section{Early Stopping Chose the Wrong Checkpoints}
\label{sec:disc-earlystop}

The protocol stops on validation loss and evaluates the checkpoint with the lowest
validation loss. Validation loss is not the reported metric. Macro F1 is, and the two do
not peak at the same epoch.

\begin{table}[htbp]
\centering
\caption{Macro F1 at the selected checkpoint against the best macro F1 reached during
training. The final column is what checkpoint selection on validation loss cost.}
\label{tab:checkpoint-cost}
\small
\begin{tabular}{@{}lrrrrr@{}}
\toprule
\textbf{Exp} & \textbf{Best val-loss epoch} & \textbf{Macro F1 there} &
\textbf{Best macro F1 epoch} & \textbf{Best macro F1} & \textbf{Cost} \\
\midrule
E1 & 28 & 0.744 & 28 & 0.744 & 0.000 \\
E2 & 36 & 0.640 & 29 & 0.672 & 0.032 \\
E3 & 24 & 0.613 & 30 & 0.650 & 0.038 \\
E4 & 46 & 0.511 & 49 & 0.571 & \textbf{0.059} \\
E5 & 20 & 0.721 & 30 & 0.724 & 0.003 \\
E6 &  6 & 0.534 & 11 & 0.601 & \textbf{0.067} \\
\bottomrule
\end{tabular}
\end{table}

E6 is the clearest case. Its validation loss reached its minimum at epoch 6 and never
improved again, so patience expired at epoch 16 and the epoch-6 weights were the ones
evaluated on test. Validation macro F1 at epoch 6 was 0.534. By epoch 11 it had reached
0.601 and it was still above 0.59 at the point training stopped. The reported test
results for E6, including its melanoma recall of 0.157, come from a model taken ten
epochs before its best measured configuration.

This is a limitation of the protocol rather than of ResNet-50, but it does not fall
evenly. The cost is 0.000 and 0.003 for two of the three MobileNetV2 runs and 0.059 and
0.067 for two of the three ResNet-50 runs. The reason is the instability of
Section~\ref{sec:disc-instability}: a noisy validation loss curve has an argmin that
lands somewhere arbitrary, whereas a smooth one has an argmin near the genuine optimum.

Part of the gap this thesis reports between the two architectures is therefore
attributable to the checkpoint selection rule interacting with optimisation noise, not to
the architectures themselves. The direction of the effect favours MobileNetV2, which is
the model the thesis concludes for, so this is a caveat that cuts against its own
conclusion and is stated for that reason.

The protocol was not changed after this was discovered. Changing the selection rule
mid-study, after seeing which runs it disadvantaged, would be a considerably worse
methodological error than the one it corrects. Section~\ref{sec:future} records
selection on validation macro F1 as the first item of future work.

\section{Why the Imbalance Remedies Hurt ResNet-50}
\label{sec:disc-remedies}

RQ2 produced an asymmetric answer. Weighted cross-entropy improved MobileNetV2's clinical
recall criterion from 0.589 to 0.648, while both remedies reduced ResNet-50's: from 0.516
to 0.378 under weighted cross-entropy and to 0.352 under oversampling. Interventions
designed to help rare classes made one architecture measurably worse at those classes.

The explanation follows from the two preceding sections. Both remedies increase gradient
variance. Weighted cross-entropy multiplies the loss on a dermatofibroma example by 12.578
against 0.214 for a nevus, a ratio of 58.8 (Table~\ref{tab:weights}), so a single rare
example can dominate a batch. Oversampling achieves a similar effect by repetition. On a
network already optimising unstably at the shared learning rate, added gradient variance
compounds the problem rather than being absorbed by it.

E4 is the endpoint of that interaction: ResNet-50's sensitivity meets the most aggressive
reweighting, it never converges within 50 epochs, and it produces the weakest result in
the study. Its basal cell carcinoma recall of 0.208 is the lowest figure in the entire
factorial, on a class it was explicitly reweighted to favour.

The general claim this supports is narrow and worth stating precisely: imbalance
remediation is not architecture-neutral. Its effect depends on how the underlying
optimisation behaves, and a remedy that helps one backbone can harm another under
otherwise identical conditions. Gap~2 of Section~\ref{sec:gaps} asked whether the best
strategy depends on the architecture. It does, and the mechanism is optimisation
stability rather than anything about the classes themselves.

\section{Precision, Recall, and What Should Be Deployed}
\label{sec:disc-tradeoff}

Table~\ref{tab:clinical} shows configurations distributed along a trade-off curve rather
than ordered from worse to better. E3 finds 63.5\% of melanomas and is correct on 34.0\%
of its melanoma calls. E6 is correct on 66.7\% of its melanoma calls and finds 15.7\% of
melanomas. On 115 test melanomas that is 73 found against 18.

Neither is unambiguously better in the abstract, which is exactly why the criterion was
fixed in advance rather than chosen once the numbers existed. Under a screening workflow
where flagged lesions receive specialist review, a false negative is an untreated cancer
and a false positive is a consultation. E3's 42 missed melanomas against E6's 97 is the
comparison that matters, and it is not close.

The honest qualification is that this reasoning depends on the downstream workflow. A
tool whose outputs are acted on without review, or one deployed where referral capacity
is the binding constraint, would weight precision differently. The criterion encodes an
assumption about deployment context, and it should be re-derived rather than inherited if
that context changes.

\section{Comparison with Published Work}
\label{sec:disc-comparison}

The accuracy figures here sit below much of the HAM10000 literature. Gao et
al.~\cite{gao2024skinlitenet} report 98.97\%; Le et al.~\cite{le2020classweightedfocal}
report 93\% averaged across classes. The best accuracy in this study is 0.801.

Direct comparison is not meaningful, for reasons established in
Section~\ref{sec:lit-ham}. Reported HAM10000 accuracy depends heavily on whether the
split is at lesion or image level, and studies frequently do not state which. An
image-level split places photographs of the same lesion on both sides of the boundary,
and on a dataset where 1{,}956 of 7{,}470 lesions are photographed more than once, that
inflates the result by an amount no reader can estimate from the paper. The lesion-level
protocol used here is the stricter choice and produces lower, better-founded numbers.

The comparison this thesis can make is internal, which is the point of running everything
through one pipeline. All twelve configurations share a split, an augmentation pipeline,
an optimiser, a schedule and a GPU. Differences between them are attributable to the
factors that varied, and that is a claim none of the cross-paper comparisons in the
surveyed corpus can make.

\section{What the Ablation Adds, and What It Does Not}
\label{sec:disc-ablation}

EfficientNet-B3 produced the highest-scoring run in the study. E11 reaches macro F1 0.738
against the best confirmatory run's 0.681, and it does so while holding 4.80 times
MobileNetV2's parameters. A thesis arguing for the lightweight model has to deal with
that honestly rather than around it.

Three things make it less awkward than it first appears.

The first is that at matched input resolution, B3 and MobileNetV2 are not statistically
distinguishable. All three of those comparisons fail to reject after Holm correction, two
of them at $p = 1.000$. B3's headline advantage comes from E11, which is a 300-pixel run,
and the resolution pairs show that the 300-pixel input is what carries it: E8 to E11 is
the single largest and only significant resolution effect in the ablation. The
architecture is not doing the work; the extra pixels are.

The second is that E11's advantage is bought with exactly the resource the thesis is
about. At 300 pixels B3 costs 8.9 times MobileNetV2's inference time on ARM64 and 4.8
times its disk footprint. For a screening tool that runs on a phone, that is not a
favourable trade at any accuracy. The question this thesis asks is what a deployable
architecture gives up, and B3 answers a different question well.

The third is a caution about the ablation's own status. E7 to E12 were added after the
confirmatory design was fixed and run. They are analysed in a separate family for that
reason (Section~\ref{sec:stat-families}), and no amount of good performance changes the
fact that they were not pre-registered. Treating E11 as the thesis's headline result
would mean promoting a post-hoc addition over the design that was specified in advance,
which is the precise failure mode the two-family structure exists to prevent.

What the ablation does establish is worth stating plainly. A third architecture, of
intermediate size and from a different design family, lands in the same place as the
other two: indistinguishable from a much smaller network at matched input size. Three
architectures spanning 2.2 to 23.5 million parameters perform comparably on this dataset,
which is stronger evidence for the capacity argument of
Section~\ref{sec:disc-headline} than the two-architecture comparison could provide alone.

It also produces the one clean methodological result in the ablation: input resolution
matters for B3 only under weighted cross-entropy, the highest-variance training condition
in the study. That is a finding the published literature cannot report, because it varies
architecture and resolution together.

\section{Limitations}
\label{sec:disc-limitations}

\paragraph{Checkpoint selection.} Documented in
Section~\ref{sec:disc-earlystop}. It disadvantages the noisier architecture and
contributes an unknown fraction of the reported architecture gap. This is the most
consequential limitation in the thesis.

\paragraph{Single seed.} Every configuration was trained once. Run-to-run variance under
a different seed is therefore unmeasured, and small differences between configurations
cannot be distinguished from seed noise. The McNemar tests address whether two
\emph{specific trained models} differ, which is a narrower claim than whether two
configurations differ in expectation. Multiple seeds per cell would cost twelve times the
compute and would be the single most valuable addition to this design.

\paragraph{E4 did not converge.} It exhausted the 50-epoch budget with validation loss
still improving. Its numbers describe an unconverged model and are reported as such.

\paragraph{Unmatched dropout.} MobileNetV2 carries $p = 0.2$, EfficientNet-B3 $p = 0.3$
and ResNet-50 none (Section~\ref{sec:method-dropout}). The asymmetry favours the models
with dropout, one of which the thesis concludes for.

\paragraph{No weight decay or schedule.} Uniform across runs and therefore not a
confound between conditions, but the probable cause of the instability that shaped
several results.

\paragraph{Rare-class sample sizes.} Dermatofibroma has 13 test images and vascular
lesions 15. A single prediction moves dermatofibroma recall by 7.7 percentage points.
Per-class figures for these classes are noisy and should not carry argumentative weight.

\paragraph{Single dataset.} HAM10000 only, by design (Gap~1). Nothing here establishes
that the ranking transfers to another dataset or another acquisition protocol.

\paragraph{Demographic coverage.} HAM10000 carries no Fitzpatrick skin-type annotation
and is drawn from two predominantly light-skinned populations. Whether these results hold
across skin tones is untested and untestable on this data. For a tool intended to widen
access to screening, this is a serious gap rather than a technical footnote.

\paragraph{No qualitative analysis.} No Grad-CAM or misclassification review was
performed, so there is no evidence about whether the models attend to lesions rather
than to acquisition artefacts such as hair, rulers or ink markings.
